\section{Identification}
\label{sec:identification}

\subsection{The singular value test}
\label{sec:svdtest}

The aging direction $v^*$ is well-defined only if the calibration interventions
agree on a common direction in methylation space. If accelerating and geroprotective
interventions push methylation in unrelated directions across contexts, the
optimisation in Equation~\eqref{eq:agingdir} has no stable solution and the
one-dimensional trajectory assumption fails.

We assess this empirically via the singular value decomposition of $A$. Let
$\lambda_1 \geq \lambda_2 \geq \cdots$ denote the singular values of $A$. The
first singular value $\lambda_1$ measures the total signed agreement across all
intervention-context pairs in the direction of $v^*$. The ratio
\begin{equation}
  \rho = \frac{\lambda_1}{\lambda_2}
  \label{eq:svdratio}
\end{equation}
is the primary test statistic for one-dimensionality. A large ratio ($\rho \gg 1$)
indicates that the aging signal is well-approximated by a single direction and the
trajectory assumption is justified. A ratio near 1 indicates that two or more
directions capture comparable amounts of signed agreement, implying that aging is
not one-dimensional in methylation space and $\gamma$ should be generalised to a
surface or higher-dimensional manifold.

In the latter case, $s(x)$ becomes a vector rather than a scalar and the
classification definitions in Section~\ref{sec:classification} require modification.
We report the full eigenvalue spectrum in Section~\ref{sec:results} and discuss the
dimensionality finding there. We note that a failed one-dimensionality test is not
a failure of the framework; it is a substantive empirical finding about the geometry
of the aging process.

\subsection{The circularity residue}
\label{sec:circularity}

The framework breaks the circularity between clocks and interventions identified
in Section~\ref{sec:related}: the aging direction $v^*$ is defined by calibration
anchors grounded in health outcome evidence, not by clock output. However, it does
not eliminate circularity entirely. A residual circularity remains: the calibration
anchors in $\mathcal{I}^+$ and $\mathcal{I}^-$ are labeled based on their
association with mortality and healthspan, and biological age is ultimately anchored
to the same outcomes.

We argue this residual circularity is unavoidable and appropriate. Any definition
of biological age that has no connection to health outcomes would be biologically
uninteresting: the concept is motivated by the observation that individuals of the
same chronological age differ in their proximity to death and disease. The question
is not whether to anchor biological age to health outcomes, but whether to do so
explicitly and directly or implicitly through a clock training procedure. Our
framework does the former.

What is eliminated is clock-dependence. The aging direction $v^*$ does not depend
on which clocks exist, how they were trained, or which CpG sites they use. Two
researchers using different clocks but the same calibration anchors and the same
cross-context identification criterion will obtain the same $v^*$ up to sampling
variability. This is a stronger guarantee than the existing literature provides.

\subsection{Local versus global coherence}
\label{sec:coherence}

The aging direction $v^*(x)$ is defined locally --- it can in principle vary across
the methylation manifold. For the integral curve $\gamma$ to be globally coherent,
$v^*$ must not rotate substantially between the methylation states characteristic
of young individuals and those characteristic of old individuals.

We test this directly by stratifying the calibration intervention datasets by
baseline age of participants. Let $v^*_{\text{young}}$ denote the aging direction
estimated from interventions applied to participants in the lowest chronological
age tertile, and $v^*_{\text{old}}$ the direction from the highest tertile. The
angle between them,
\begin{equation}
  \theta = \arccos\!\left(
    \frac{\langle v^*_{\text{young}},\, v^*_{\text{old}} \rangle_c}
         {|v^*_{\text{young}}|_c \cdot |v^*_{\text{old}}|_c}
  \right)
  \label{eq:curvature}
\end{equation}
measures the curvature of the trajectory between young and old methylation states.
A small $\theta$ supports the use of a single globally estimated $v^*$; a large
$\theta$ indicates that the aging direction rotates substantially with age and that
a locally varying vector field is required.

We report $\theta$ in Section~\ref{sec:results}. Regardless of the outcome, the
principal curve estimator in Section~\ref{sec:trajectory} accommodates local
variation by construction: it follows the data rather than a fixed linear direction.
The curvature test therefore bears primarily on the interpretability of $v^*$ as a
single global direction rather than on the validity of the trajectory estimate.

\subsection{Geometry of the state space}
\label{sec:geometry}

The weighted inner product in Equation~\eqref{eq:inner} provides a Euclidean
approximation to the geometry of $\mathcal{M}$. This approximation is reasonable
for CpG sites with beta values in the interior of $[0,1]$, where small displacements
are approximately linear. It breaks down near the boundaries: a site with mean
beta value $0.02$ cannot decrease by more than $0.02$, so variance-normalised
displacements at such sites are not comparable to those at interior sites.

We address this in the empirical analysis by excluding CpG sites with mean beta
below $0.05$ or above $0.95$ in the control samples of any context. A more principled
treatment would replace the Euclidean geometry with a Riemannian metric appropriate
for the beta distribution, such as the Fisher information metric on the simplex or
the arc-sine transformation. We note this as a natural refinement and leave the
Riemannian extension for future work.

A separate geometric question is whether the methylation manifold is well-approximated
by a smooth submanifold at all, or whether it has a more complex topology (branching,
holes, or disconnected components corresponding to different cell types). Cell-type
heterogeneity in bulk methylation samples is a known confounder, and we address it
in the empirical pipeline via standard cell-type deconvolution before computing
displacements. After deconvolution, the residual manifold structure is expected to
be smooth enough for the principal curve estimator to apply reliably.
